% ===== FILE 6/9: Động lực SADGS, Kim tự tháp đa tỉ lệ, Structure tensor gộp, Trị riêng & bước sóng, Chiếu trục Gaussian =====

% --- Slide: Động lực SADGS ---
\begin{frame}{SADGS: bổ sung tín hiệu tần số cho Adaptive Density Control}
\scriptsize
Nhắc lại hạn chế đã nêu ở cuối phần trước: gradient một mình không đo được "Gaussian này đang to/nhỏ \textbf{so với chi
tiết ảnh cần biểu diễn} theo tỉ lệ nào". \textbf{Ý tưởng cốt lõi của SADGS} (Structure-Aware Densification): đo trực
tiếp một tỉ số hình học giữa \textbf{kích thước Gaussian sau khi chiếu lên màn hình} và \textbf{kích thước chi tiết ảnh
cục bộ} (một dạng "bước sóng" texture) tại đúng vị trí Gaussian đó --- gọi tỉ số này là $\eta$.
\begin{itemize}\itemsep2pt
    \item Nếu Gaussian \textbf{to hơn} chi tiết cần biểu diễn ($\eta>1$): nó đang "nuốt chửng", làm mờ mất chi tiết đó
    --- cần \textbf{split} thành các mảnh nhỏ hơn, và quan trọng hơn: chia \textbf{đúng theo trục} bị vi phạm, không
    phải chia đều mọi trục như cơ chế gốc.
    \item Nếu Gaussian \textbf{nhỏ hơn nhiều} so với chi tiết cục bộ ($\eta\le0.1$): nó đang "phục vụ thừa" một vùng gần
    như không đổi --- một Gaussian lớn hơn cũng đủ sức biểu diễn vùng đó --- là ứng viên \textbf{prune}.
    \item \textbf{Điểm quan trọng cần nhấn mạnh:} SADGS \textbf{không thay thế} khung ADC gốc vừa học --- nó \textbf{kế
    thừa toàn bộ} khung clone/split/prune, chỉ \textbf{bổ sung thêm} một điều kiện kích hoạt dựa trên $\eta$, kết hợp
    bằng phép \texttt{OR} với điều kiện gradient cũ (chi tiết đầy đủ ở phần sau).
    \item Bốn slide tiếp theo xây dựng $\eta$ từ con số $0$: ảnh gốc $\to$ kim tự tháp đa tỉ lệ $\to$ structure tensor
    $\to$ trị riêng $\to$ bước sóng cục bộ $\to$ chiếu trục Gaussian lên màn hình $\to$ $\eta$.
\end{itemize}
\end{frame}

% --- Slide: Kim tự tháp đa tỉ lệ ---
\begin{frame}{Bước 1: kim tự tháp đa tỉ lệ (multiscale) trên ảnh GT}
\scriptsize
Trước khi có thể so sánh "kích thước Gaussian" với "kích thước chi tiết ảnh", cần một cách đo \textbf{chi tiết ảnh có
kích thước bao nhiêu} tại mỗi vị trí. Hàm \texttt{get\_multiscale\_structure\_tensor\_v1}
(\texttt{utils/loss\_utils.py:232--}) làm mờ dần ảnh gốc qua nhiều tầng $i=0,\dots,\texttt{levels}-1$ (mặc định $3$
tầng), tầng sau mờ hơn tầng trước theo một tỉ lệ cố định:
\[
\sigma_i=\sigma_0\cdot(\texttt{octave\_step})^i,\quad \texttt{octave\_step}=1.5,\qquad
b_i = \big\lVert I^{(i-1)}_{\text{mờ}} - I^{(i)}_{\text{mờ}}\big\rVert_2 \ \ (\text{DoG, xấp xỉ Laplacian})
\]
\begin{itemize}\itemsep1pt
    \item $b_i$ (\emph{Difference-of-Gaussians}, kỹ thuật kinh điển trong thị giác máy tính, ví dụ dùng trong SIFT): là
    \textbf{bản đồ năng lượng cạnh} tại tầng $i$ --- vùng phẳng cho $b_i\approx0$, vùng có cạnh/chi tiết rõ cho $b_i$ lớn.
    \item Tần số ứng với mỗi tầng: $f_i=1/(\texttt{octave\_step})^i$ --- tầng $i$ càng lớn (mờ càng nhiều) thì $f_i$
    càng nhỏ, ứng với chi tiết \textbf{thô} hơn.
    \item \textbf{Vì sao cần nhiều tỉ lệ chứ không chỉ một?} Một hoạ tiết lặp (ví dụ vân gỗ, gạch lát) có thể "rõ" khi
    nhìn ở độ phân giải mịn nhưng gần như "biến mất" khi nhìn ở độ phân giải thô, hoặc ngược lại với cấu trúc lớn (viền
    một toà nhà). Gộp nhiều tầng tránh bỏ sót chi tiết chỉ tồn tại rõ ràng ở một quy mô riêng.
\end{itemize}
\begin{center}
\includegraphics[width=0.4\linewidth]{01_pyramid_levels.png}
\end{center}
\end{frame}

% --- Slide: Structure tensor gộp đa tỉ lệ ---
\begin{frame}{Bước 2: gộp thành một structure tensor duy nhất $\hat J$}
\scriptsize
Tại mỗi tầng $i$, ta có thể tính một \textbf{structure tensor} $J_i=\begin{psmallmatrix}S_{xx}&S_{xy}\\S_{xy}&S_{yy}\end{psmallmatrix}_i$
--- một ma trận $2\times2$ cổ điển trong xử lý ảnh, đo \textbf{hướng} và \textbf{cường độ} biến thiên cục bộ (dựng từ
tích ngoài của vector gradient ảnh, làm mờ nhẹ để ổn định). Structure tensor gộp toàn bộ các tầng theo trọng số
(\texttt{loss\_utils.py:283--300}):
\[
\boxed{\;\hat J = \frac{\sum_i \big(J_i/\mathrm{tr}J_i\big)\,w_i\,f_i^2}{\sum_i w_i}\;},\qquad w_i = b_i^{\,p}\ (p=3\text{, mặc định})
\]
\begin{itemize}\itemsep1pt
    \item $J_i/\mathrm{tr}J_i$: \textbf{chuẩn hoá theo vết} --- mỗi tầng chỉ còn đóng góp thông tin \emph{hướng/hình
    dạng} (dị hướng của biến thiên), loại bỏ ảnh hưởng của việc tầng nào đó có năng lượng tuyệt đối lớn hơn tầng khác một cách tình cờ.
    \item $w_i=b_i^3$: tầng nào có cạnh/chi tiết rõ ràng hơn ($b_i$ lớn hơn) được \textbf{"tin cậy" nhiều hơn} khi gộp;
    luỹ thừa $3$ khuếch đại sự chênh lệch này khá mạnh, để tầng "nhiễu" (ít cấu trúc) không làm loãng thông tin từ tầng
    có cấu trúc rõ.
    \item $f_i^2$: hệ số quy đổi để các tầng --- vốn bị làm mờ ở các $\sigma_i$ khác nhau --- được đưa về \textbf{cùng
    một thang tần số vật lý} trước khi cộng gộp.
    \item Kết quả $\hat J=(S_{xx},S_{xy},S_{yy})$ chính là "nguyên liệu" dùng lại xuyên suốt các bước tính $\eta$ tiếp theo.
\end{itemize}
\begin{center}
\includegraphics[width=0.4\linewidth]{01_structure_tensor_channels.png}
\end{center}
\end{frame}

% --- Slide: Trị riêng & bước sóng cục bộ ---
\begin{frame}{Bước 3: trị riêng của $\hat J$ $\to$ bước sóng cục bộ $w_{\min}$}
\scriptsize
Từ ma trận $\hat J$, trích ra \textbf{trị riêng lớn nhất} $\lambda_1$ (\texttt{utils/freq\_utils.py:313--334}) --- đại
lượng vô hướng duy nhất cần cho bước tiếp theo:
\[
\lambda_1=\frac{\mathrm{tr}\hat J}{2}+\sqrt{\Big(\frac{\mathrm{tr}\hat J}{2}\Big)^2-\det\hat J},\qquad
\boxed{\;w_{\min}=\dfrac{1}{\sqrt{\lambda_1}+10^{-5}}\;}
\]
\begin{itemize}\itemsep1pt
    \item $\lambda_1$ tỉ lệ với \textbf{năng lượng gradient bình phương} theo hướng biến thiên mạnh nhất tại điểm đó ---
    vùng có cạnh sắc, chi tiết dày đặc sẽ cho $\lambda_1$ lớn; vùng phẳng cho $\lambda_1\approx0$.
    \item Phép biến đổi "năng lượng $\to$ bước sóng" dùng nghịch đảo căn bậc hai: trực giác là năng lượng gradient cao
    tương ứng với \textbf{tần số không gian cao} (chi tiết mịn, đổi nhanh theo pixel), mà tần số cao nghĩa là \textbf{bước
    sóng ngắn} --- do đó $\lambda_1$ lớn $\Rightarrow$ $w_{\min}$ \textbf{nhỏ}. Ngược lại, vùng phẳng $\lambda_1\to0$ cho
    $w_{\min}\to10^5$ (rất lớn --- "không có chi tiết nào cần bám theo ở đây cả").
    \item $w_{\min}$ sẽ đóng vai trò \textbf{mẫu số}, là "cây thước" để so sánh với kích thước Gaussian ở bước tính $\eta$ ngay sau đây.
\end{itemize}
\begin{center}
\includegraphics[width=0.4\linewidth]{02_wavelength_eigen.png}
\end{center}
\end{frame}

% --- Slide: Chiếu trục Gaussian lên màn hình ---
\begin{frame}{Bước 4: chiếu 3 trục chính của Gaussian lên màn hình}
\scriptsize
Có "cây thước" $w_{\min}$ rồi, giờ cần đo "kích thước Gaussian" theo đúng đơn vị tương thích (pixel trên màn hình).
Hàm \texttt{compute\_projected\_axes\_subset} (\texttt{utils/freq\_utils.py:116--178}) dùng lại đúng Jacobian $J$ của
phép chiếu pinhole đã học ở phần đầu bài giảng:
\[
\boxed{\;a_j = J\,(R_{\text{view}}R_i)\,\mathrm{diag}(s)\,e_j\;},\qquad \ell_j = \lVert a_j\rVert = \sqrt{u_j^2+v_j^2+10^{-8}}
\]
\begin{itemize}\itemsep1pt
    \item $R_i\,\mathrm{diag}(s)\,e_j$: véc-tơ trục chính thứ $j$ của ellipsoid ($j=1,2,3$, ba trục vuông góc của
    ellipsoid trong không gian world), có độ dài $=s_j$.
    \item $R_{\text{view}}$ xoay véc-tơ đó sang hệ toạ độ camera; $J$ tuyến tính hoá phép chiếu phối cảnh, đánh giá tại
    đúng vị trí $\mu$ của Gaussian --- giống hệt cách dùng $J$ trong EWA splatting ở phần render.
    \item $\ell_j$ là \textbf{độ dài tính bằng pixel}, sau khi chiếu, của trục chính thứ $j$ --- đây chính là tử số
    của $\eta$.
    \item Mỗi Gaussian có \textbf{3 giá trị} $\ell_1,\ell_2,\ell_3$ (ứng với 3 trục chính hình học, \textbf{không phải}
    3 kênh màu R/G/B) --- nhờ đó SADGS biết được Gaussian đang vi phạm \emph{theo hướng cụ thể nào}, để sau này split
    \emph{đúng trục} đó thay vì chia đều mọi hướng.
\end{itemize}
\begin{center}
\includegraphics[width=0.4\linewidth]{02_jacobian_axes.png}
\end{center}
\end{frame}
